\def\year{2018}\relax
%File: formatting-instruction.tex
\documentclass[letterpaper]{article} %DO NOT CHANGE THIS
\usepackage{aaai18}  %Required
\usepackage{times}  %Required
\usepackage{helvet}  %Required
\usepackage{courier}  %Required
\usepackage{url}  %Required
\usepackage{graphicx}  %Required
\usepackage{amsfonts}
\usepackage{multirow}
\usepackage{multicol}
\usepackage{amssymb}
\usepackage{amsmath}
\usepackage{makecell}
\frenchspacing  %Required
\setlength{\pdfpagewidth}{8.5in}  %Required
\setlength{\pdfpageheight}{11in}  %Required
%PDF Info Is Required:
  \pdfinfo{
/Title (2018 Formatting Instructions for Authors Using LaTeX)
/Author (AAAI Press Staff)}
\setcounter{secnumdepth}{0}  
 \begin{document}
% !TEX spellcheck = English
% The file aaai.sty is the style file for AAAI Press 
% proceedings, working notes, and technical reports.
%
\title{Path-Based Attention Neural Model for Fine-Grained Entity Typing}
\author{Denghui Zhang$^{1}$, Manling Li$^{1}$, Pengshan Cai$^{2}$, Yantao Jia$^{1}$, Yuanzhuo Wang$^{1}$, Xueqi Cheng$^{1}$\\
$^{1}$Institute of Computing Technology, Chinese Academy of Sciences, Beijing, 100190, China\\
$^{2}$School of Computer Science, University of Massachusetts Amherst, MA 01003\\
\{zhangdenghui, wangyuanzhuo, cxq\}@ict.ac.cn, \{limanlingcs, jamaths.h\}@gmail.com, pengshancai@cs.umass.edu\\
}
\maketitle
\begin{abstract}
Fine-grained entity typing aims to assign entity mentions in the free text with types arranged in a hierarchical structure. 
%It suffers from the label noise in training data generated by distant supervision. 
%Traditional distant supervision based methods employ a structured data source as a weak supervision and do not need hand-labeled data, but they neglect the label noise in the automatically labeled training corpus. 
Traditional distant supervision based methods employ a structured data source as a weak supervision and do not need hand-labeled data, but they neglect the label noise in the automatically labeled training corpus. 
Although recent studies use many features to prune wrong data ahead of training, they suffer from error propagation and bring much complexity. 
%Although recent studies use many features to reduce the noise, they suffer from the error propagation phenomenon. 
%Although recent studies use attention to reduce the noise, they neglect the hierarchical structure of types. 
%Although recent studies manage to overcome it by pruning wrong data ahead of training. suffer from error propagation and bring much parameter and computation complexity.
%In other distant supervision based tasks, e.g., relation extraction, attention neural models are used to reduce noise through an end-to-end process, however, they neglect the hierarchical structure of types, and can not be easily applied into fine-grained entity typing. %which is distinct in fine-grained entity typing. 
In this paper, we propose an end-to-end typing model, called the path-based attention neural model (PAN), to learn a noise-robust performance by leveraging the hierarchical structure of types. Experiments %on two data sets 
demonstrate its effectiveness. 
%Fine-grained entity typing aims to label entity mentions with types according to their context, and the type set constitutes a multi-layered hierarchy. It suffers from the noise in training data generated by distant supervision, and previous works manage to overcome it by pruning wrong data ahead of training. However, these approaches lead to error propagation and bring much parameter and computation complexity. In this paper, we propose a path-based attention neural model (PAN), which is end-to-end and robust to noise. Experiments on two data sets demonstrate its validity.
\end{abstract}

\section{Introduction}
\noindent Fine-grained entity typing aims to assign types (e.g., ``person'', ``politician'', etc.) %(e.g., ``person'', ``businessman'') 
to entity mentions in the local context (a single sentence), and the type set constitutes a tree-structured hierarchy (i.e., type hierarchy). %tree-structured hierarchy. 
Recent years witness the boost of neural models in this task, e.g., \cite{Shimaoka2016Neural} employs an attention based LSTM to attain sentence representations and achieves state-of-the-art performance. However, it still suffers from noise in training data,  which is a main challenge in this task. %impediment
%Although the application of neural models achieve great success, e.g., \cite{Shimaoka2016Neural} employs an attention mechanism to attain sentence representations and achieves the state-of-the-art performance, it is still far from satisfactory. 
%Lacking in manually annotated data, training sets are automatically generated by distant supervision, i.e., labeling entity mentions with type labels in existing knowledge bases. However, this approach introduces label noise to entity mentions, since the type labels obtained from knowledge bases may be irrelevant to the local context. 
The training data is generated by distant supervision, which assumes that if an entity has a type in knowledge bases (KBs), then all sentences containing this entity will express this type. %then all mentions of this entity in different sentences will express this type. 
%, and thus label entity mentions with types in existing KBs.
This method inevitably introduces irrelevant types to the context. 
For example, the entity ``Donald Trump'' has types ``person", ``businessman'' and ``politician'' in KBs, thus all three types are annotated for its mentions in the training corpora. But in sentence ``Donald Trump announced his candidacy for President of US.'', only ``person'' and ``politician'' are correct types, % for the mention ``Donald Trump'',
while ``businessman'' can not be deduced from the sentence, thus serves as noise. 
%Trump started his career at his father's real estate development company, 
To alleviate this issue, a few systems try to denoise training data by filtering irrelevant types ahead of training. %adding preprocessing mechanisms. 
%Among these methods, \cite{Ren2016Label} achieves the state-of-the-art performance, 
%The typical method is \cite{Ren2016Label}, which 
For instance, \cite{Ren2016Label} proposes PLE to  identify correct types by jointly embedding mentions, context and type hierarchy, and then use clean data to train classifiers. %an embedding method to identify correct types for each mention, then use clean data to train classifiers. %identify correct types for each mention, and then clean training data is generated . %, and proven to be state-of-the-art in reducing noise. 
%of each training instance. %embeds mentions, types and features into the same vector space to identify the correct types of each training instance. 
%This approach brings improvements, but it split the whole task into two steps: denoise preprocessing and training. Fine tuning the parameters of the two disjunct parts can be a difficult and time-consuming process.
However, the denoising and training process are not unified, which may cause error propagation and bring much additional complexity.
%detached  seperate disjunct
%this drawback correct?

Motivated by this, we propose an end-to-end typing model, called the path-based attention neural model (PAN), %by leveraging the hierarchical structure of types to learn a noise-robust performance. 
to select relevant sentences to each type, which can dynamically reduce the weights of wrong labeled sentences for each type during training. 
This idea is inspired by some successful attempts to reduce noise in relation extraction, e.g.,\cite{Lin2016Neural}. %proposes an attention mechanism to focus on relevant sentences to a single relation. 
However, these methods fail to formulate type hierarchy, which is distinct in fine-grained entity typing. 
%Although this idea has been successful used in relation extraction to reduce noise, e.g.,\cite{Lin2016Neural}, the
Specifically, %the types are hierarchical dependent, %i.e., the existence of a subtype indicates the existence of its parent type. 
%a sentence indicating a subtype will also express its parent type. 
if a sentence indicates a type, its parent type can be also deduced from the sentence. 
Like the example above, ``politician'' is the subtype of ``person''. Since the sentence indicates that ``Donald Trump'' is ``politician'',  ``person'' should also be assigned. 
%i.e., an entity mention labeled with a subtype will also be labeled with its parent type. %Like the example above, %the mention ``Donald Trump'' indicates two types ``person'' and ``politician'', and 
%For example, ``person'' is the parent type of ``politician''. If subtype ``politician'' is labeled, the parent ``person'' must be also labeled. %Simple application will lose the hierarchical information.
%Considering this hierarchical constraint, path-based attention for each type is obtained based on its path to the top coarse type (e.g., ``person'', ``businessman'') in type hierarchy, which enables parameter sharing of types in the same path, and thus reduce noise effectively and yields a better typing classifier.  
Thus, we build path-based attention for each type by utilizing its path to its coarsest parent type (e.g., person, businessman) in the type hierarchy. Compared to the simple attention in relation extraction, %it can capture the hierarchical information of types and 
it enables parameter sharing for types in the same path. With the support of hierarchical information of types, %distinct type hierarchy in fine-grained entity typing, 
it can reduce noise effectively and yields a better typing classifier. 
Experiments on two data sets validate effectiveness of PAN. 
%Actually, there are some successful attempts to reduce noise by an end-to-end neural model in relation extraction, e.g.,\cite{Lin2016Neural} proposes an attention mechanism to focus on relevant sentences to a single relation. However, these methods fails to formulate hierarchy among types, which is distinct in fine-grained entity typing. %can not be easily applied to fine-grained entity typing with two problems: 
%In detail, the types are hierarchical dependent, i.e., the existence of a subtype indicates the existence of its parent type. Like the example above, the mention ``Donald Trump'' indicates two types ``person'' and ``politician'', and ``person'' is the parent type of ``politician''. If subtype ``politician'' is labeled, the parent ``person'' must be also labeled. %Simple application will lose the hierarchical information.

%Motivated by this, we propose a path-based attention neural model (PAN) to focus on relevant sentences to type paths, % to make full use of the hierarchical dependencies of types, 
%and explore the representation of type paths to impose hierarchical dependency constraints on types. %restrict types in the same path to be bounded to each other during training. %types in the same path is necessarily bounded to each other during training. %by transforming the multi-label problem to multi-path problem, % by employing type hierarchy in knowledge bases, 
%During training process, it can dynamically reduce the weights of wrong labeled sentences, and the attention learned for each type could utilize the hierarchical structure to enable parameter sharing. Thus it can reduce noise through an end-to-end process. %restrict types in the same path to be bounded to each other. % without preprocessing.
%Experiments validate effectiveness of PAN. %over two data sets validate its effectiveness. %the effectiveness of PAN.


\section{Path-Based Attention Neural Model}
The architecture of PAN is illustrated in \figurename \ref{nn}. 
Supposing that there are $n$ sentences containing entity $e$, i.e., $\mathcal{S}_e=\{s_1,s_2,...,s_n\}$, and $\mathcal{T}_e$ is the automatically labeled types based on KBs. 
%As illustrated in \figurename \ref{nn}, 
Firstly PAN employs LSTM to generate representations of sentences $\mathbf{s}_i$ following \cite{Shimaoka2016Neural}, where  $\mathbf{s}_i \in \mathbb{R}^{d}$ is the semantic representation of $s_i$, $i \in \{1,2,...,n\}$. 
%PAN firstly employs LSTM to generate  representations of sentences, as illustrated in \figurename \ref{nn}. 
Afterwards, we build path-based attention $\alpha_{i,t}$ over sentences $s_i$ for each type $t  \in \mathcal{T}_e$, which is expected to focus on relevant sentences to type $t$. %de-emphasize the noisy sentence irrelevant to type $t$. 
Then, the representation of sentence set $\mathcal{S}_e$ for type $t$, denoted by $\mathbf{s}_{e,t} \in \mathbb{R}^{d}$, 
%for sentence set $\mathcal{S}_e$, its representation $\mathbf{s}_{e,t} \in \mathbb{R}^{d}$ 
is calculated through weighted sum of vectors of sentences. %these sentence vectors. 
Finally, we obtain predicted types through a classification layer. 
%In this paper, we propose a sentence-level attention-based convolutional neural network (CNN) for distant supervised relation extraction. As illustrated in Fig. 1, we employ a CNN to embed the semantics of sentences. Afterwards, to utilize all informative sentences, we represent the relation as semantic composition of sentence em- beddings. To address the wrong labelling prob- lem, we build sentence-level attention over multiple instances, which is expected to dynamically reduce the weights of those noisy instances. Fi- nally, we extract relation with the relation vector weighted by sentence-level attention. 
%Suppose there are $n$ sentences in certain sentence for entity $e$ in training data, denoted by $C=\{c_1,c_2,...,c_n\}$. %i.e., instead of using all of them to train a classifer, we first pack them into a bag $S=\{s_1,s_2,...,s_n\}$, then use only the representation of bag $S$ to train. 
\vspace{-10pt}
\begin{figure}[!h]
\centering
\scalebox{0.55}{\includegraphics{nn.pdf}}
\vspace{-12pt}
\caption{The architecture of PAN for given entity $e$, type $t$}
\label{nn}
\vspace{-10pt}
\end{figure} 

%The idea of PAN is that different path should focus on distinct sentences of an entity, 
%\subsection{Path-Based Attention}
%In detail, given $\mathcal{D}_e=\{(m_i,c_i,\mathcal{P}_e)\}$, $i \in \{1,...,n\}$, for each path $p \in \mathcal{P}_e$, 
More precisely, given $e$, an attention $\alpha_{i,t}$ is learned to score how well sentence $s_i$ matches type $t$, i.e., %focus on relevant sentences $s_i$ to each type $t$. 
\[
\alpha_{i,t} = \frac{\exp(\mathbf{s}_i\mathbf{A}\mathbf{p}_t)}{\sum_{j=1}^{n} \exp(\mathbf{s}_j\mathbf{A}\mathbf{p}_t)},
\vspace{-4pt}
\]
%where $\mathbf{s}_i \in \mathbb{R}^{d}$ is generated through LSTM following \cite{Shimaoka2016Neural}, 
where $\mathbf{A} \in \mathbb{R}^{d \times d}$ is a weighted diagonal matrix. 
$\mathbf{p}_t \in \mathbb{R}^{d}$ is the representation of path $p_t$ for type $t$. 
Specifically, for each type, we define one path as a sequence of types starting from its coarsest parent type, and ending with it. %e.g., $person$, $location$, etc. 
More formally, for type $t_l$, $p_{t_l}=t_1$$\rightarrow$$t_2$$\rightarrow$$...$$\rightarrow$$t_l$, where $t_1$ is its coarsest parent type, and $t_{i+1}$ is the subtype of $t_{i}$. 
%Note that there is exactly one path for each type, as types are arranged in a tree-structured hierarchy. 
For example, for type $t_l=politician$, its path is $p_{t_l}=person$$\rightarrow$$politician$.
%Due to the hierarchical dependency between types, the sentences indicating subtypes will help to predict the parent type, and thus these sentences should also be focused on by the parent type. 
%the sentences should be focused on by a parent type if the subtypes focuses on. 
%Namely, the learning of attention parameters of a subtype restricts its parent type to perform the same learning, 
%the parameters of parent type must also be bounded to updates of child types.% when update the attention of child types. 
%so the types in the same path are necessarily bounded to each other during training process. 
%It can be seen that a path is composed of the entity types. 
%Therefore 
We represent the path $p_{t_l}$ as a semantic composition of all the types on the path, i.e., $\mathbf{p}_{t_l}=\mathbf{t}_1 \circ \mathbf{t}_2 \circ...\circ \mathbf{t}_l$, 
%\vspace{-8pt}
%\[
%\mathbf{p}=\mathbf{t}_1 \circ \mathbf{t}_2 \circ...\circ \mathbf{t}_l, 
%\vspace{-2pt}
%\]
where $\mathbf{t}_i \in \mathbb{R}^{d}$ is the representation of type $t_i$, which is a parameter to learn. $\circ$ is a  composition operator. 
%In this paper, we consider two operations Addition (PAN-A) and Multiplication (PAN-M).
In this paper, we consider two  types of operators: (1) Addition (PAN-A), where $\mathbf{p}_{t_l}$ equals the sum of type vectors. %i.e., $\mathbf{p}_{t_l}=\mathbf{t}_1+\mathbf{t}_2+...+\mathbf{t}_l$; 
(2) Multiplication (PAN-M), where $\mathbf{p}_{t_l}$ equals the cumulative product of type vectors. %i.e., $\mathbf{p}_{t_l}=\mathbf{t}_1 \cdot \mathbf{t}_2 \cdot...\cdot \mathbf{t}_l$. %More operations are not discussed here due to length of the paper. 
In this way, path-based attention enables the model to share parameters between types in the same path. For example, the attention learned for ``person'' could assist the learning of the attention for ``politician''. It makes learning easier especially for infrequent subtypes, which suffer from dearth of training data, since the attentions for these subtypes can get support from the attention for parent type. 
% draw support from
%the attention learned for parent types could assist the learning of attention for subtypes.  
%For example, if ``person'' focus on a set of sentences $S$, then the sentences, which ``politician'' focuses on, should lies in $S$. 

Then, the representation of sentence set $\mathcal{S}_e$ for type $t$, i.e., $\mathbf{s}_{e,t}$, is calculated through weighted sum of sentence vectors,
%the vectors of sentences, namely, 
\vspace{-6pt}
%\[
%\mathbf{c}^{(p)}=\sum_{i=0}^{n} \alpha_{i}^{(p)}\mathbf{c}_i, \quad \alpha_{i}^{(p)} = \frac{\exp(\mathbf{c}_i\mathbf{A}\mathbf{p})}{\sum_{j=0}^{n} \exp(\mathbf{c}_j\mathbf{A}\mathbf{p})},
%\]
\[
\mathbf{s}_{e,t}=\sum_{i=1}^{n} \alpha_{i,t}\mathbf{s}_i. %\quad \alpha_{i,t} = \frac{\exp(\mathbf{c}_i\mathbf{A}\mathbf{p}_t)}{\sum_{j=1}^{n} \exp(\mathbf{c}_j\mathbf{A}\mathbf{p}_t)},
\vspace{-6pt}
\]

Since one mention can have multiple types, we employ a classification layer consisting of $N$ logistic classifiers, where $N$ is the total number of types. Each classifier outputs the probability of respective type, i.e., 
\vspace{-4pt}
%\[
%P(p|\mathbf{c}^{(p)})=\frac{\exp(\mathbf{w}_{p}^\mathrm{T} \mathbf{c}^{(p)}+\mathbf{b}_{p})}{1+\exp(\mathbf{w}_{p}^\mathrm{T} \mathbf{c}^{(p)}+\mathbf{b}_{p})},
%\]
\[
P(t|\mathbf{s}_{e,t})=\frac{\exp(\mathbf{w}_{t}^\mathrm{T} \mathbf{s}_{e,t}+\mathbf{b}_{t})}{1+\exp(\mathbf{w}_{t}^\mathrm{T} \mathbf{s}_{e,t}+\mathbf{b}_{t})},
\vspace{-4pt}
\]
where $\mathbf{w}_{t}, \mathbf{b}_{t}  \in \mathbb{R}^{d}$ are the logistic regression parameters. % of path $p$, $\mathbf{s}$ is the attention sum of a sentence bag defined previously.
%Since one sentence may indicate multiple paths, we define multi-path loss function by cross entropy as follows,
To optimize the model, a multi-type loss is defined according to the cross entropy as follows,
\vspace{-4pt}
\[
%J=-\sum_{e \in \mathcal{E}} \sum_{p \in \mathcal{P}_e} [\mathbb{I}_{p}\ln P(p|\mathbf{c}^{(p)})+(1-\mathbb{I}_{p})\ln(1-P(p|\mathbf{c}^{(p)}))],
%J=-\sum_{e} \sum_{p} [\mathbb{I}_{p}\ln P(p|\mathbf{c}^{(p)})+(1-\mathbb{I}_{p})\ln(1-P(p|\mathbf{c}^{(p)}))],
J=-\sum_{e} \sum_{t} [\mathbb{I}_{t}\ln P(t|\mathbf{s}_{e,t})+(1-\mathbb{I}_{t})\ln(1-P(t|\mathbf{s}_{e,t}))],
\vspace{-4pt}
\]
%$\mathcal{E}$ is the set of bags, 
where $\mathbb{I}_t$ is indicator function to indicate whether $t$ is the annotated type of entity $e$, i.e., $t\in \mathcal{T}_e$. 

\section{Experiments and Conclusion}
Experiments are carried on two widely used datasets OntoNotes and FIGER(GOLD), and the training dataset of OntoNotes is noisy compared to FIGER(GOLD) \cite{Shimaoka2016Neural}. The statistics of the datasets are listed in \tablename \ref{dataset}. %provided by \cite{Shimaoka2016Neural}. 
%Note that the training dataset of OntoNotes is noisy compared to FIGER (GOLD). %(since FIGER (GOLD) uses anchors to detect entities whereas OntoNotes uses an external tool)
%Note that OntoNotes is more sensitive to noise, since it is carefully manually annotated according to the context.   %FIGER(GOLD) has large training data labeled by a 2-level type hierarchy, while the training data of OntoNotes is smaller and labeled by a 3-level hierarchy. %Moreover, the manually annotated test data in OntoNotes is more context dependent than FIGER(GOLD) and thus more noise-sensitive.
%  a contains 200,000 Freebase entities that were mapped to 102 types (2-level hierarchy), and a set of 4,300,000 sentences containing those entities. OntoNotes is smaller dataset containing 223,342 entities that were mapped to 89 types (3-level hierarchy), and a set of 143,709 sentences. OntoNotes is carefully manually annotated according to context, which is more context dependent.  
%During training, we use ``raw'' noisy data provided by \cite{Shimaoka2016Neural}. %we use the automatically induced publicly available datasets ("raw" noisy data) provided by \cite{Shimaoka2016Neural}
\vspace{-25pt}
\begin{table}[h!]
\small
\centering
%\setlength{\belowcaptionskip}{-10pt}
\caption{Statistics of the datasets.}
%\begin{tabular}{|c|c|c|c|c|c|}
\begin{tabular}{p{1.9cm}|p{0.75cm}p{0.75cm}p{0.83cm}p{0.78cm}p{0.78cm}}
\hline
Datasets & \#Type &\#Layer & \#Context & \#Train & \#Test  \\
\hline
OntoNotes & 89 & 3& 143K & 223K & 8,963  \\
FIGER(GOLD) &  113 & 2& 1.51M & 2.69M & 563  \\
\hline
\end{tabular}
\label{dataset}
\vspace{-14pt}
\end{table}

We employ Strict Accuracy (Acc), Loose Macro F1 (Ma-F1), and Loose Micro F1 (Mi-F1) as evaluation measures following \cite{Shimaoka2016Neural}. Specifically, ``Strict" evaluates on the type set of each entity mention, while ``Loose'' on each type. ``Marco'' is the geometric average over all mentions, while ``Micro'' is the arithmetic average. %``Strict" is computed over whole type set, while ``Loose'' over single type. Ma-F1 is the average over the averages of each mention, and Mi-F1 is the average over all mentions.
%only considers prediction true when the whole predicted type set is exactly same as ground truth, while ``Loose'' considers each predicted type, and Ma-F1 is computed as the average of each mention, Mi-F1 is computed as the average over all mentions.
%\begin{itemize}
%\item Strict: The prediction is considered correct if and only if the predicted label set is the same as the true label set.
%\item Loose Macro: The Macro-Precision (Ma-P) and MacroRecall (Ma-R) are computed for each mention, then their F1 is computed.
%\item Loose Macro: The Micro-Precision (Mi-P) and MicroRecall (Mi-R) scores are averages over all mentions, then their F1 is computed.
%\end{itemize}
The baselines are chosen from two aspects: (1) Predicting types in a unified process using raw noisy data, i.e., TLSTM \cite{Shimaoka2016Neural}, and other methods shown in \tablename \ref{performance}. 
%To demonstrate the superiority of path-based attention, 
(2) Predicting types using clean data by denoising ahead, i.e., H\_PLE and F\_PLE \cite{Ren2016Label}. To prove the superiority of path-based attention, we also directly apply the attention neural model in relation extraction \cite{Lin2016Neural} without using type hierarchy (AN). %During experiment, $d=300$, and the hidden-size of the LSTM is 100. %We used Adam as our optimization method with a learning rate of 0.001, a mini-batch size of 1,000. %and hidden-size of the LSTM is 100, we get the parameters by evaluating on the validation set.
The results of baselines are the best results reported in their papers. %\cite{Shimaoka2016Neural} and \cite{Ren2016Label}. % following their papers. 
%We compare our model with three baseline methods: 1)Figer: Using only hand-crafted features to train a perception classifier 2) LSTM \cite{Shimaoka2016Neural}: Combine hand-crafted features with sentence representation learned by an Attentive LSTM network 3)PLE+Figer \cite{Ren2016Label}: Use Partial Label Embedding to produce denoised training data, then use the data to train a perception classifier.
\vspace{-17pt}
\begin{table}[htbp]
\small
\centering
%\setlength{\belowcaptionskip}{-10pt}
\caption{Performance on FIGER(GOLD) and OntoNotes}
%\vspace{-7pt}
\label{performance}
%\vspace{-4pt}
 \begin{tabular}{p{1cm}|p{0.61cm}p{0.87cm}p{0.87cm}|p{0.60cm}p{0.87cm}p{0.87cm}}
\hline
\multirow{2}{*}{Metric} & \multicolumn{3}{c|}{OntoNotes} & \multicolumn{3}{c}{FIGER(GOLD)} \\
%\multirow{2}{*}{Metric} & \multicolumn{3}{c}{OntoNotes} & \multicolumn{3}{|c|}{FIGER(GOLD)} \\
%\cline{2-7}
  & Acc &  Ma-F1 & Mi-F1 & Acc & Ma-F1 & Mi-F1 \\
\hline
HYENA & 24.9 & 49.7 & 44.6 & 28.8 & 52.8 & 50.6 \\
FIGER & 36.9 & 57.8 & 51.6 &  47.4 & 69.2 & 65.5  \\
%Figer & 36.90 & 57.80 & 51.60 &  52.3 & 69.9 & 69.3  \\
TLSTM & 50.8 & 70.1 & 64.9 &  59.7 & 79.0 & 75.4  \\
\hline
AN & 52.3 & 71.7 & 65.2 & 60.0 & 79.5 & 75.9  \\
%PAN-A & \textbf{54.5} & \textbf{72.7} & \textbf{66.3} & \textbf{59.8} & \textbf{79.2} & \textbf{76.1}  \\
PAN-A & \textbf{54.9} & \textbf{72.8} & \textbf{66.5} & \textbf{60.2} & \textbf{79.9} & \textbf{76.2}  \\
PAN-M &  53.0 & 71.9 &65.3 & 60.0 & 79.4 &  76.0 \\
\hline
\end{tabular}
\vspace{-10pt}
\end{table}

We can observed that: (1) When using the same raw noisy data, PAN outperforms all methods on both data sets, %It is unsurprising since PAN is robust to noise
which proves the anti-noise ability of PAN. 
(2) PAN performs better than AN, since the attention learned in PAN utilizes the hierarchical structure to enable parameter sharing. %PAN considers hierarchical dependency among types while training. %path-based attention restrict types in the same 
(3) The improvements on OntoNotes are higher than FIGER(GOLD), %it can also be easily understood 
because OntoNotes is more noisy, and the hierarchical structure in OntoNotes is more complex with more layers, which further demonstrates that path-based attention does well with type hierarchy, and proves the superiority of PAN in reducing noise. 
(4) PAN-A achieves better performance than PAN-M, which shows that addition operator can better capture type hierarchy. 


\vspace{-17pt}
\begin{table}[htbp]
\small
\centering
%\setlength{\belowcaptionskip}{-10pt}
\caption{Performance on FIGER(GOLD) and OntoNotes}
%\vspace{-7pt}
\label{performancedenoise}
%\vspace{-4pt}
 \begin{tabular}{p{1cm}|p{0.61cm}p{0.87cm}p{0.87cm}|p{0.60cm}p{0.87cm}p{0.87cm}}
\hline
\multirow{2}{*}{Metric} & \multicolumn{3}{c|}{OntoNotes} & \multicolumn{3}{c}{FIGER(GOLD)} \\
%\multirow{2}{*}{Metric} & \multicolumn{3}{c}{OntoNotes} & \multicolumn{3}{|c|}{FIGER(GOLD)} \\
%\cline{2-7}
  & Acc &  Ma-F1 & Mi-F1 & Acc & Ma-F1 & Mi-F1 \\
\hline
H\_PLE & 54.6 & 69.2 & 62.5 & 54.3 & 69.5 & 68.1 \\
F\_PLE & \textbf{57.2} & 71.5 & 66.1& 59.9 & 76.3 & 74.9 \\
\hline
PAN-A & 54.9 & \textbf{72.8} & \textbf{66.5} & \textbf{60.2} & \textbf{79.9} & \textbf{76.2}  \\
%\Xhline{1.5pt}
%PAN-A & 54.5 & \textbf{72.1} & \textbf{66.2} & 59.80 & \textbf{79.17} & \textbf{76.02}  \\
%PAN-M &  &  &  &  &  &   \\
\hline
\end{tabular}
\vspace{-10pt}
\end{table}

%Compared with the performance on the denoised data by PLE, 
%Compared with H\_PLE and F\_PLE using denoised data to train, 
As shown in \tablename  \ref{performancedenoise}, %PAN outperforms H\_PLE and F\_PLE on Ma-F1 and Mi-F1, though PAN uses raw noisy data and H\_PLE and F\_PLE uses denoised data. 
PAN using raw noisy data outperforms H\_PLE and F\_PLE using denoised data on Ma-F1 and Mi-F1. 
%achieves comparable performance by using raw noisy data with H\_PLE and F\_PLE using denoised training data, and PAN even has higher  Ma-F1 and Mi-F1. %on Ma-F1 and Mi-F1, PAN even has better performance. %However, the end-to-end process makes PAN Considering that PAN reduce noise through an end to end process, PAN is more 
It makes sense that F\_PLE has higher Acc on OntoNotes since the noise is reduced before training, but it needs to learn additional parameters about mentions, context and types, %apart from training process,  
while PAN only needs to learn parameters of attention. Thus, PAN  is more efficient to reduce noise. %employ clean data to train. %since PLE reduces the data noise before training, while PAN try to denoise during training through an end-to-end process. 
%considering the efficiency, PAN is more effective to denoise. %In conclusion, PAN can reduce the noise in training data effectively through an end-to-end process.
%On OntoNotes, using the same raw noisy training data, adding sentence-level attention mechanism to an Attentive LSTM network outperforms the Attentive LSTM network on all the three criteria. Although without any denoise preprosessing, we outperform PLE+Figer on Loose Macro and Loose Micro, and we achieve comparable results on Strict Accuracy. This confirms the denoise ability of our model.  2) The improvements on FIGER(GOLD) are not salient, we speculate the reason is that OntoNote is noisy compared to FIGER (GOLD) (since FIGER (GOLD) uses anchors to detect entities whereas OntoNotes uses an external tool), even the attention LSTM outperforms PLE+Figer.
%Furthermore, our model can be trained end-to-end and does not need any extra complex preprocessing, and the sentence-level attention mechanism does not bring much extra complexity. Thus, it is more suitable for large-scale application.
% Specifically, it improves the Strict Accuracy by 3 percent, Loose Macro 2 percent,  Loose Micro 1 percent.
% Furthermore, we perform extensive analysis of the attention mechanism of our model and establish that the attention mechanism learns to focus on the decisive sentences, which is highly relevant to the certain type path and the mention in it can be labeled by  distinct type obviously. %Actually, the small bag containing %both which are known to be highly relevant to successfully classifying a mention.

In conclusion, PAN can reduce noise effectively through an end-to-end process, and achieves better typing performance on datasets with more noise. %or more complex type hierarchy. 

% \section{Acknowledgments}
% The work was funded by National Natural Science Foundation of China (No. 61572469, 61402442, 91646120, 61572473, 61402022), the National Key R\&D Program of China (No. 2016QY02D0405, 2016YFB1000902), and National Grand Fundamental Research 973 Program of China (No. 2013CB329602, 2014CB340401).

%section{Conclusion}
%In this paper, we propose a sentence-level attention mechanism to establish a soft alignment that enables the model to focus on informative and decisive sentence according to type paths, which can effectively reduce the noise in an end-to-end process and fulfill the hierarchical dependency constraint of types automatically. Experiments demonstrate the effectiveness of it.

%References and End of Paper
%These lines must be placed at the end of your paper
\bibliography{reference}
\bibliographystyle{aaai}

\end{document}
